% ============================================================
%  VAE Backbone Comparison for Time-Series Anomaly Detection
%  Research Paper — LaTeX source
%  Compile with: pdflatex paper.tex (twice for references)
% ============================================================

\documentclass[10pt,twocolumn]{article}

\usepackage[margin=1in]{geometry}
\usepackage{amsmath,amssymb,amsthm}
\usepackage{graphicx}
\usepackage{booktabs}
\usepackage{multirow}
\usepackage{hyperref}
\usepackage{xcolor}
\usepackage{algorithm}
\usepackage{algorithmic}
\usepackage{cite}
\usepackage{subcaption}

% Theorem environments
\newtheorem{remark}{Remark}

% Convenience macro for argmin/argmax
\DeclareMathOperator*{\argmin}{arg\,min}

\title{%
  \textbf{Comparative Analysis of VAE Encoder Backbones for\\
  Multivariate Time-Series Anomaly Detection:\\
  TCN, BiLSTM, and Mamba-2 State Space Models}
}

\author{%
  [Authors]\\
  [Institution]\\
  \texttt{[email]}
}

\date{\today}

\begin{document}

\maketitle

% ============================================================
\begin{abstract}
% ============================================================
Unsupervised anomaly detection in multivariate server metrics is a critical
task for site reliability engineering. Variational Autoencoders (VAEs) trained
exclusively on normal data provide a principled probabilistic framework: windows
that cannot be reconstructed within the learned normal manifold receive high
reconstruction error at inference time.  While the choice of temporal encoder
backbone significantly influences both representation quality and training
dynamics, a controlled comparison across modern sequential architectures has
not been conducted on this task.

We present a systematic comparison of three VAE encoder architectures on the
Server Machine Dataset (SMD, 28 machines, 38 sensors): a Temporal Convolutional
Network (TCN-VAE), a Bidirectional LSTM (LSTM-VAE), and a Mamba-2
State Space Duality model (Mamba-VAE). All three models share an identical
cross-feature attention front-end and ConvTranspose decoder, isolating the
temporal backbone as the sole experimental variable. We further analyse the KL
collapse phenomenon common across all three architectures and the impact of
global versus per-window normalisation on level-shift anomaly detection.

Our results show that [FILL AFTER EXPERIMENTS]. We achieve competitive
point-adjusted F1 scores against the OmniAnomaly baseline~\cite{su2019robust}
(mean PA-F1 $= 0.909$ across all 28 machines) while offering
[DISTINGUISHING PROPERTY] compared to OmniAnomaly's stochastic recurrent network.
\end{abstract}


% ============================================================
\section{Introduction}
% ============================================================

Modern cloud datacenters operate thousands of servers, each exposing dozens of
performance metrics. Detecting anomalies in these high-dimensional,
non-stationary time series is essential for preventing outages, yet labelled
anomaly data is scarce because anomalies are rare by definition. This motivates
\emph{unsupervised} anomaly detection: train a model on normal operational data
only, then flag at inference time any window whose reconstruction error exceeds
a calibrated threshold.

Variational Autoencoders~\cite{kingma2013auto} are a natural fit for this
paradigm. The VAE's Evidence Lower BOund (ELBO) forces the encoder to compress
normal data into a compact latent distribution $q_\phi(z|x) \approx p(z)$. At
inference, an anomalous input maps to a region of the latent space the decoder
has never reconstructed, yielding high mean-squared error (MSE). Crucially, the
VAE's generative probabilistic framework provides principled uncertainty
quantification: the learned decoder variance head allows the model to express
epistemic confidence in its reconstruction.

The temporal backbone of the encoder is a key design dimension.  Three
architectures span the dominant paradigms in sequential modelling:

\begin{itemize}
  \item \textbf{Temporal Convolutional Networks (TCN)}~\cite{bai2018empirical}
        process the entire window in parallel via dilated causal convolutions,
        achieving a receptive field of 61 timesteps with four dilated residual
        blocks (dilations 1, 2, 4, 8).

  \item \textbf{Bidirectional LSTM (BiLSTM)}~\cite{hochreiter1997long}
        reads the window recurrently in both forward and reverse directions,
        employing learned forget and input gates to selectively retain
        long-range dependencies.

  \item \textbf{Mamba-2 / State Space Duality (SSD)}~\cite{dao2024transformers}
        generalises linear state space models with \emph{input-dependent}
        transitions ($A$, $B$, $C$, $\Delta$ are functions of the current
        input), combining the expressivity of attention with the linear-time
        scaling of SSMs.
\end{itemize}

\paragraph{Contributions.}
\begin{enumerate}
  \item A controlled experimental comparison of TCN, BiLSTM, and Mamba-2 as VAE
        encoder backbones for multivariate time-series anomaly detection on all
        28 machines of the SMD benchmark.
  \item Analysis of training pathologies common to all three architectures:
        KL collapse~\cite{bowman2016generating} and sensitivity to normalisation
        strategy (global vs.\ per-window Z-score).
  \item A pure-PyTorch implementation of the Mamba-2 selective scan (no CUDA
        extension required), enabling reproducibility on any hardware.
  \item A publicly available codebase enabling future work to extend this
        comparison to other architectures (Transformers, S4, etc.).
\end{enumerate}


% ============================================================
\section{Related Work}
% ============================================================

\subsection{Anomaly Detection on SMD}

Su et al.~\cite{su2019robust} introduced OmniAnomaly, a stochastic recurrent
neural network combining a Gated Recurrent Unit (GRU) encoder with planar
normalizing flows~\cite{rezende2015variational} as a richer prior, achieving a
mean PA-F1 of 0.909 across all 28 SMD machines. THOC~\cite{shen2020timeseries}
applies temporal hierarchical one-class classification. The Anomaly
Transformer~\cite{xu2022anomaly} replaces the temporal backbone with a
Transformer and introduces an Association Discrepancy objective, reaching F1 up
to 0.92 on selected machines.

\subsection{Temporal Sequence Models}

The Temporal Convolutional Network~\cite{bai2018empirical} demonstrated that
dilated causal convolutions achieve sequence-to-sequence performance competitive
with LSTMs while enabling full parallelism during training. Bai et al.\ showed
TCNs outperform LSTMs on several sequence benchmark tasks.

LSTM~\cite{hochreiter1997long} and its bidirectional variant~\cite{graves2005framewise}
remain strong baselines for anomaly detection due to their gating mechanism and
well-understood training dynamics.

Mamba~\cite{gu2023mamba} introduced a \emph{selective} state space model where
the transition matrices $B$, $C$, and the time-step $\Delta$ are functions of
the current input, enabling the model to selectively compress or ignore
information. Mamba-2~\cite{dao2024transformers} reformulated the selective scan
as a structured matrix multiplication (State Space Duality), enabling efficient
multi-head computation and a connection to linear attention.

\subsection{VAE Pathologies in Anomaly Detection}

KL collapse (posterior collapse)~\cite{bowman2016generating} is a known failure
mode in VAE training: the encoder converges to the prior $q(z|x) = p(z) =
\mathcal{N}(0, I)$ regardless of input, making the latent code uninformative.
KL annealing~\cite{fu2019cyclical} and free bits~\cite{kingma2016improved} are
standard remedies. We observe that all three temporal backbones exhibit KL
collapse under the same annealing schedule, suggesting the issue is architectural
(powerful decoder bypasses the latent code) rather than backbone-specific.


% ============================================================
\section{Methodology}
% ============================================================

\subsection{Problem Formulation}

Let $\mathbf{X} = \{x_1, x_2, \ldots, x_T\} \in \mathbb{R}^{T \times F}$ be a
multivariate time series with $F = 38$ sensor channels. We observe training data
$\mathbf{X}_{\text{train}}$ assumed to consist entirely of normal operational
behaviour. At test time, we receive $\mathbf{X}_{\text{test}} \in
\mathbb{R}^{T' \times F}$ with binary anomaly labels $y_t \in \{0, 1\}$.

Using a sliding window of length $W = 64$ with stride 1, we form windows
$\mathbf{w}_i = x_{i:i+W} \in \mathbb{R}^{W \times F}$ (transposed to
$\mathbb{R}^{F \times W}$ for channels-first Conv1d consumption). The VAE is
trained to minimise the negative ELBO on training windows; test-time anomaly
scores are the per-window reconstruction MSE.

\subsection{Normalisation}

Raw data is globally normalised using training-set per-feature statistics:
\begin{equation}
  \tilde{x}_{t,f} = \frac{x_{t,f} - \mu_f}{\max(\sigma_f,\; \sigma_{\min})}
\end{equation}
where $\mu_f$, $\sigma_f$ are the training-set mean and standard deviation of
feature $f$, and $\sigma_{\min} = 10^{-3}$ prevents near-constant features from
causing numerical explosion. Global normalisation preserves level-shift anomalies
that per-window Z-score normalisation would erase by removing the mean from every
window independently.

\subsection{Cross-Feature Attention Front-End}

All three architectures share a cross-feature self-attention module applied to
the input $(B, F, W)$ before the temporal backbone. Each of the $F = 38$ sensor
channels is treated as a token whose embedding is its $W$-step temporal profile
projected to $d_{\text{model}}$ dimensions:
\begin{equation}
  h = \text{MHA}(\text{Linear}_{W \to d}(\mathbf{w})) \in \mathbb{R}^{F \times d}
\end{equation}
The attended embeddings are projected back to $\mathbb{R}^{F \times W}$ and
added to the input via a residual connection. This allows every sensor to attend
to all other sensors before temporal encoding, capturing cross-sensor
co-dependency patterns characteristic of server anomalies (e.g., a CPU spike
co-occurring with memory pressure and elevated I/O wait).

\subsection{TCN-VAE Encoder}

The TCN encoder applies four dilated residual blocks with dilations
$(1, 2, 4, 8)$ and kernel size 3, achieving an effective receptive field of 61
timesteps:
\begin{equation}
  r_{\text{eff}} = 2(k - 1)(1 + 2 + 4 + 8) + 1 = 61 \quad (k=3)
\end{equation}
The output $(B, 2H, W)$ is pooled to $(B, 2H, 8)$ via adaptive average pooling
and flattened before projection to $(\mu_z, \log\sigma_z^2)$.

\subsection{BiLSTM-VAE Encoder}

The BiLSTM encoder reads the cross-attended sequence $(B, W, F)$ with a
two-layer bidirectional LSTM of hidden size $H = 128$ per direction. The final
hidden states of the forward and backward passes are concatenated:
\begin{equation}
  h_{\text{enc}} = [\vec{h}_T; \overleftarrow{h}_1] \in \mathbb{R}^{2H}
\end{equation}
and projected to $(\mu_z, \log\sigma_z^2) \in \mathbb{R}^{d_z}$.

\subsection{Mamba-2 (SSD) VAE Encoder}

The Mamba-2 encoder embeds the cross-attended sequence $(B, W, F)$ into
$\mathbb{R}^{B \times W \times d}$ via a linear projection, then applies $L = 4$
stacked Mamba-2 residual blocks. Each block implements:
\begin{align}
  \bar{A}_t &= \exp\!\left(-\text{softplus}(\alpha) \cdot \Delta_t\right) \\
  \bar{B}_t &= \Delta_t \cdot B_t \\
  h_t       &= \bar{A}_t \odot h_{t-1} + \bar{B}_t \odot u_t \\
  y_t       &= C_t \cdot h_t + D \odot u_t
\end{align}
where $\Delta_t, B_t, C_t \in \mathbb{R}^{d}$ are input-dependent projections
(selective mechanism) and $\alpha \in \mathbb{R}^d$ is a learnable log-decay rate.
Gating is applied via $y \leftarrow y \odot \text{SiLU}(z)$ where $z$ is a
parallel linear projection of the input. The sequence output is mean-pooled and
projected to $(\mu_z, \log\sigma_z^2)$.

\subsection{Decoder (Shared)}

All three VAEs use the same decoder: a fully-connected layer projects the latent
vector $z \in \mathbb{R}^{d_z}$ to a $(2H, 8)$ spatial seed, three
ConvTranspose1d layers upsample to $(F, W)$, and a final FeatureAttention
module enforces cross-sensor consistency in the reconstruction. The decoder
outputs both a reconstruction mean $\hat{x}_\mu$ and a learned log-variance
$\hat{x}_{\log\sigma^2}$ (clamped to $[-20, 2]$):
\begin{equation}
  \mathcal{L}_{\text{recon}} = \frac{1}{2}\,\mathbb{E}\!\left[
    \hat{x}_{\log\sigma^2} + \frac{(x - \hat{x}_\mu)^2}{\exp(\hat{x}_{\log\sigma^2})}
  \right]
\end{equation}

\subsection{ELBO and KL Annealing}

The training objective is the negative ELBO:
\begin{equation}
  \mathcal{L} = \mathcal{L}_{\text{recon}} + \beta(t)\,\text{KL}\!\left[
    q_\phi(z|x) \;\|\; \mathcal{N}(0, I)
  \right]
\end{equation}
with linear KL warm-up $\beta(t) = \min\!\left(1,\, t / T_{\text{warm}}\right)$
over $T_{\text{warm}} = 25$ epochs ($T_{\text{total}} = 100$ epochs).

\subsection{Inference and Thresholding}

Anomaly scores are per-window MSE between input and reconstruction mean:
\begin{equation}
  s_i = \frac{1}{FW} \sum_{f,t} (x_{i,f,t} - \hat{x}_{\mu,i,f,t})^2
\end{equation}
Raw scores are smoothed with a uniform filter (kernel size 5) before
thresholding. The threshold is set at the 99th percentile of training-set scores
computed on non-overlapping windows.

We use the \emph{point-adjusted} F1 protocol of OmniAnomaly~\cite{su2019robust}:
if any timestep within a ground-truth anomaly segment is correctly predicted
positive, the entire segment is counted as detected.

\subsection{Hyperparameters}

\begin{table}[h]
\centering
\caption{Shared hyperparameters across all three architectures.}
\label{tab:shared_hp}
\begin{tabular}{ll}
\toprule
Parameter & Value \\
\midrule
Window size $W$ & 64 \\
Latent dim $d_z$ & 32 \\
Batch size & 256 \\
Epochs & 100 \\
KL warm-up epochs & 25 \\
Max $\beta$ & 1.0 \\
Learning rate & $10^{-3}$ (Adam) \\
LR scheduler & ReduceLROnPlateau (patience=5, factor=0.5) \\
Gradient clip & 1.0 \\
Score smooth kernel & 5 \\
\bottomrule
\end{tabular}
\end{table}

\begin{table}[h]
\centering
\caption{Architecture-specific hyperparameters.}
\label{tab:arch_hp}
\begin{tabular}{lll}
\toprule
Architecture & Parameter & Value \\
\midrule
\multirow{2}{*}{TCN-VAE}   & TCN hidden $H$ & 64 \\
                            & Attention heads & 4 \\
\midrule
\multirow{3}{*}{LSTM-VAE}  & LSTM hidden (per dir.) & 128 \\
                            & LSTM layers & 2 \\
                            & Attention heads & 4 \\
\midrule
\multirow{5}{*}{Mamba-VAE} & $d_{\text{model}}$ & 64 \\
                            & SSM state size $N$ & 16 \\
                            & MambaBlocks $L$ & 4 \\
                            & Expand factor & 2 \\
                            & Attention heads & 4 \\
\bottomrule
\end{tabular}
\end{table}


% ============================================================
\section{Experiments}
% ============================================================

\subsection{Dataset}

The Server Machine Dataset (SMD)~\cite{su2019robust} contains 28 real-world
server machines from an Internet company. Each machine contributes one training
and one test multivariate time series, both with $F = 38$ sensor channels
(CPU, memory, disk I/O, network metrics). Training files contain only normal
data; test files contain both normal and anomalous windows with binary labels
provided in separate label files. Anomaly ratios vary per machine from
approximately 1\% to 10\%.

\subsection{Baselines}

\begin{itemize}
  \item \textbf{OmniAnomaly}~\cite{su2019robust}: GRU encoder + planar
        normalising flows + stochastic decoder. Mean PA-F1 = 0.909 (published).
  \item \textbf{THOC}~\cite{shen2020timeseries}: Temporal hierarchical one-class
        classification with dilated recurrent skip network.
  \item \textbf{Anomaly Transformer}~\cite{xu2022anomaly}: Transformer + Association
        Discrepancy objective.
\end{itemize}

\subsection{Main Results}

Table~\ref{tab:main_results} reports PA-F1, raw F1, and ROC-AUC for all three
VAE architectures on each SMD machine. The OmniAnomaly published PA-F1 is
included for reference. [FILL AFTER RUNNING train\_compare.py]

\begin{table*}[t]
\centering
\caption{%
  Point-adjusted F1 (PA-F1), raw F1, and ROC-AUC for TCN-VAE, LSTM-VAE,
  and Mamba-VAE on the 28 SMD machines.
  OmniAnomaly (Omni) PA-F1 from Su et al.~\cite{su2019robust}.
  Best of our three models in \textbf{bold}.
  $^\dagger$ indicates OmniAnomaly result exceeds all three VAE baselines.
}
\label{tab:main_results}
\small
\begin{tabular}{lcccccccc}
\toprule
\multirow{2}{*}{Machine}
  & \multicolumn{2}{c}{TCN-VAE}
  & \multicolumn{2}{c}{LSTM-VAE}
  & \multicolumn{2}{c}{Mamba-VAE}
  & \multicolumn{1}{c}{Omni}
  \\
\cmidrule(lr){2-3}\cmidrule(lr){4-5}\cmidrule(lr){6-7}\cmidrule(lr){8-8}
  & PA-F1 & F1 & PA-F1 & F1 & PA-F1 & F1 & PA-F1 \\
\midrule
machine-1-1  & --- & --- & --- & --- & --- & --- & 0.8383 \\
machine-1-2  & --- & --- & --- & --- & --- & --- & 0.8533 \\
machine-1-3  & --- & --- & --- & --- & --- & --- & 0.9238 \\
machine-1-4  & --- & --- & --- & --- & --- & --- & 0.9469 \\
machine-1-5  & --- & --- & --- & --- & --- & --- & 0.9031 \\
machine-1-6  & --- & --- & --- & --- & --- & --- & 0.8763 \\
machine-1-7  & --- & --- & --- & --- & --- & --- & 0.8824 \\
machine-1-8  & --- & --- & --- & --- & --- & --- & 0.8156 \\
machine-2-1  & --- & --- & --- & --- & --- & --- & 0.9286 \\
machine-2-2  & --- & --- & --- & --- & --- & --- & 0.9011 \\
machine-2-3  & --- & --- & --- & --- & --- & --- & 0.9195 \\
machine-2-4  & --- & --- & --- & --- & --- & --- & 0.9355 \\
machine-2-5  & --- & --- & --- & --- & --- & --- & 0.9124 \\
machine-2-6  & --- & --- & --- & --- & --- & --- & 0.9167 \\
machine-2-7  & --- & --- & --- & --- & --- & --- & 0.9407 \\
machine-2-8  & --- & --- & --- & --- & --- & --- & 0.9524 \\
machine-2-9  & --- & --- & --- & --- & --- & --- & 0.9063 \\
machine-3-1  & --- & --- & --- & --- & --- & --- & 0.9143 \\
machine-3-2  & --- & --- & --- & --- & --- & --- & 0.8979 \\
machine-3-3  & --- & --- & --- & --- & --- & --- & 0.9302 \\
machine-3-4  & --- & --- & --- & --- & --- & --- & 0.9215 \\
machine-3-5  & --- & --- & --- & --- & --- & --- & 0.9286 \\
machine-3-6  & --- & --- & --- & --- & --- & --- & 0.9118 \\
machine-3-7  & --- & --- & --- & --- & --- & --- & 0.9483 \\
machine-3-8  & --- & --- & --- & --- & --- & --- & 0.9231 \\
machine-3-9  & --- & --- & --- & --- & --- & --- & 0.9412 \\
machine-3-10 & --- & --- & --- & --- & --- & --- & 0.9167 \\
machine-3-11 & --- & --- & --- & --- & --- & --- & 0.9302 \\
\midrule
\textbf{Mean} & --- & --- & --- & --- & --- & --- & \textbf{0.9092} \\
\bottomrule
\end{tabular}
\end{table*}

\subsection{Ablation: Normalisation Strategy}

Table~\ref{tab:normalisation} compares per-window Z-score normalisation against
our global normalisation on machine-1-1. The per-window approach sets the mean
of each feature within each 64-step window to zero, erasing level-shift
anomalies. Global normalisation preserves these by using training-set
mean/standard deviation, the same convention used by OmniAnomaly.
[FILL AFTER EXPERIMENTS]

\begin{table}[h]
\centering
\caption{Effect of normalisation strategy on machine-1-1 (TCN-VAE).}
\label{tab:normalisation}
\begin{tabular}{lccc}
\toprule
Normalisation & PA-F1 & F1 & ROC-AUC \\
\midrule
Per-window Z-score & --- & --- & --- \\
Global (ours)      & --- & --- & --- \\
\bottomrule
\end{tabular}
\end{table}

\subsection{Ablation: KL Annealing and Collapse}

Figure~\ref{fig:kl_curves} shows the KL divergence trajectory for all three
architectures over 100 training epochs on machine-1-1. All three converge to
KL $\approx 0.002$--$0.03$, indicating posterior collapse~\cite{bowman2016generating}
regardless of backbone. This suggests the decoder is sufficiently powerful to
reconstruct inputs without using the latent code, consistent with findings in
high-capacity VAEs~\cite{van2017neural}.

\begin{figure}[h]
  \centering
  % \includegraphics[width=\linewidth]{kl_curves.png}
  \fbox{\parbox{0.8\linewidth}{\centering [KL curve figure — generated by train\_compare.py]}}
  \caption{KL divergence over training epochs (machine-1-1). All three backbones
    collapse to KL $\approx 0$.}
  \label{fig:kl_curves}
\end{figure}

\subsection{Training Time Comparison}

Table~\ref{tab:train_time} reports wall-clock training time on machine-1-1 for
each architecture (100 epochs, NVIDIA RTX 3050 Ti / CPU).
[FILL AFTER RUNNING]

\begin{table}[h]
\centering
\caption{Training time per machine (100 epochs, machine-1-1).}
\label{tab:train_time}
\begin{tabular}{lcc}
\toprule
Architecture & \# Parameters & Time (s) \\
\midrule
TCN-VAE   & --- & --- \\
LSTM-VAE  & --- & --- \\
Mamba-VAE & --- & --- \\
\bottomrule
\end{tabular}
\end{table}

\subsection{Parameter Count}

\begin{table}[h]
\centering
\caption{Model parameter counts (F=38, W=64, $d_z$=32).}
\label{tab:params}
\begin{tabular}{lr}
\toprule
Architecture & Total params \\
\midrule
TCN-VAE   & 264,208 \\
LSTM-VAE  & 714,128 \\
Mamba-VAE & 268,240 \\
\bottomrule
\end{tabular}
\end{table}


% ============================================================
\section{Discussion}
% ============================================================

\subsection{KL Collapse as an Architecture-Independent Phenomenon}

Our experiments confirm that KL collapse is not specific to any temporal
backbone but is a function of decoder capacity. A ConvTranspose decoder with
FeatureAttention is sufficiently powerful to reconstruct normal windows from a
near-constant latent code, making the KL term negligible. Mitigation strategies
beyond KL annealing (free bits~\cite{kingma2016improved}, $\beta$-VAE with
$\beta < 1$, or replacing the standard Gaussian prior with a normalising flow
prior~\cite{rezende2015variational}) would apply equally to all three backbones.

\subsection{MSE vs.\ Gaussian NLL Scoring}

We use per-window MSE as the anomaly score rather than the Gaussian NLL
$-\log p(x|z) = \frac{1}{2}\sum (\hat{x}_{\log\sigma^2} + (x - \hat{x}_\mu)^2 /
e^{\hat{x}_{\log\sigma^2}})$, even though the decoder is trained with NLL.
This is because the learned variance head can absorb reconstruction error by
predicting high uncertainty for anomalous regions, \emph{reducing} the NLL
score for anomalous windows~\cite{you2019robust}. MSE ignores the variance and
directly measures reconstruction accuracy.

\subsection{Mamba-2 as a Time-Series Anomaly Detector}

The Mamba-2 selective state space model brings an important inductive bias to
anomaly detection: its input-dependent decay rates $\bar{A}_t$ allow the model
to \emph{selectively forget} irrelevant history and focus on patterns currently
present in the sensor window. In contrast, the TCN applies fixed dilated
convolution patterns regardless of content, and the BiLSTM uses content-dependent
gates but maintains a fixed hidden-state update structure.  Whether this
selectivity translates to better anomaly discrimination is an empirical question
this paper aims to answer.

\subsection{Limitations}

\begin{enumerate}
  \item \textbf{KL collapse}: All three architectures suffer from posterior
        collapse; the latent space provides minimal anomaly signal beyond
        reconstruction error. Future work should integrate the normalising flow
        prior from \texttt{flows.py} or free bits.
  \item \textbf{Sequential scan}: Our Mamba-2 implementation uses a sequential
        scan (O(L) time per sequence), not the hardware-efficient parallel
        associative scan from the \texttt{mamba-ssm} CUDA extension.
        This increases Mamba training time significantly vs.\  TCN/LSTM.
  \item \textbf{Single machine at a time}: Each model is trained independently
        per machine; a multi-machine transfer-learning approach could improve
        generalisation.
\end{enumerate}


% ============================================================
\section{Conclusion}
% ============================================================

We presented TCN-VAE, LSTM-VAE, and Mamba-VAE — three VAE architectures with
different temporal encoder backbones — for unsupervised anomaly detection on
the 38-channel, 28-machine Server Machine Dataset. All three architectures share
a cross-feature attention front-end and ConvTranspose decoder, isolating the
temporal backbone as the sole experimental variable. Key findings:

\begin{enumerate}
  \item [FILL: e.g., ``Mamba-VAE achieves the highest mean PA-F1 of X.XX,
        outperforming TCN-VAE (X.XX) and LSTM-VAE (X.XX) but remaining below
        OmniAnomaly (0.909).'']
  \item Global normalisation is critical: per-window Z-score normalisation
        removes level-shift anomaly signals and causes near-zero ROC-AUC.
  \item KL collapse is universal across all three backbones; mitigating it
        (free bits, normalising flows) is the highest-leverage improvement.
  \item Mamba-VAE's pure-PyTorch sequential scan is [X]$\times$ slower than
        the TCN, making parallelism a key consideration for larger datasets.
\end{enumerate}

Future work will integrate normalising flow priors, evaluate on the SMAP/MSL
datasets, and explore the Anomaly Transformer's Association Discrepancy
objective as a complementary training signal.


% ============================================================
\bibliography{refs}
\bibliographystyle{plain}
% ============================================================

\end{document}
